\documentclass[12pt]{article}

% --- Packages ---
\usepackage[margin=1in]{geometry}
\usepackage{booktabs}
\usepackage{array}
\usepackage[table,dvipsnames]{xcolor}
\usepackage{graphicx}
\usepackage{float}
\usepackage{amsmath}
\usepackage{hyperref}

% --- Title ---
\title{Supplementary Information: Data Quality Assurance Analysis\\
\large Keheala Study 2}
\date{}

\begin{document}

\maketitle
\tableofcontents
\newpage

This supplementary document presents descriptive statistics and data quality assurance (DQA) analyses conducted as part of the Keheala Study 2 randomized controlled trial, which evaluated a digital adherence support platform for tuberculosis (TB) patients in Kenya. The DQA assesses recording errors in TIBU, Kenya's national electronic TB reporting system, by treating a contemporaneous paper registry audit as the ground truth. The TIBU dataset available for this study covers patients registered between January 2017 and December 2021 ($N = 119{,}653$ nationwide). A paper registry audit was conducted across study clinics between November 2018 and September 2019, yielding $3{,}645$ audited records, of which $3{,}370$ could be matched to a corresponding TIBU record and form the basis of the DQA. Tables~\ref{tbl:patient_char} and~\ref{tbl:outcomes} characterize the study population across three nested cohorts: all patients registered in TIBU during the study period, patients at study clinics in TIBU (the MITT sample), and the subset of study-clinic patients whose records were audited in the DQA.

The core of the DQA compares treatment outcomes recorded in TIBU against outcomes recorded in the paper registry for each matched patient. Three types of discrepancy are distinguished. Type I errors occur when TIBU records a successful outcome (treatment completion or cure) but the paper registry records a failure (death, treatment failure, or loss to follow-up), representing a false-positive in the digital record. Type II errors are the reverse: TIBU records a failure while the paper registry records success, representing a false-negative. Incomplete-data discrepancies arise when TIBU contains no outcome but the paper registry does. The cross-tabulation in Table~\ref{tbl:crosstab} maps the full distribution of these discrepancies across outcome categories; subsequent tables stratify error rates by clinic characteristics (province, urbanicity, and size) and patient characteristics (sex, age, disease type, and HIV status); and the time-series figures trace how error rates vary with the age of the patient record and across the study period.

% =============================================================================
\section{Patient Characteristics}
% =============================================================================

Tables~\ref{tbl:patient_char} and~\ref{tbl:outcomes} describe the study
population across three cohorts: all patients registered in the TIBU electronic
system during the study period, patients at study clinics recorded in TIBU
(MITT sample), and patients at study clinics whose records were audited in the
paper registry DQA.

\renewcommand{\thetable}{1a}
\begin{table}[H]
    \centering
    \caption{Individual and disease characteristics across TIBU and paper registry cohorts.}
    \label{tbl:patient_char}
    \scriptsize{
\begin{tabular}{lrr|rr|rr}
\hline\hline \\[-8pt]
Characteristic & \multicolumn{4}{c}{Digital} & \multicolumn{2}{c}{Paper} \\
 & \multicolumn{2}{c}{\shortstack{All \\ (N=119,653)}} & \multicolumn{2}{c}{\shortstack{Study clinics \\ (N=14,962)}} & \multicolumn{2}{c}{\shortstack{Study clinics \\ (N=3,263)}} \\
 & $N$ & (\%) & $N$ & (\%) & $N$ & (\%) \\ \hline \\[-8pt]
\textit{Individual characteristics} & & & & & & \\
\quad Male & 75,597 & (63.2\%) & 9,514 & (63.6\%) & 2,037 & (64.5\%) \\
\quad Mean age (years) & 36.1 & & 34.4 & & 32.7 & \\
\rule{0pt}{14pt}\textit{Disease characteristics} & & & & & & \\
\quad Bacteriologically confirmed & 62,149 & (51.9\%) & 11,243 & (75.1\%) & 2,223 & (70.4\%) \\
\quad Drug resistant & 1,343 & (1.1\%) & 2 & (0.0\%) & 3 & (0.1\%) \\
\quad Extrapulmonary TB & 17,558 & (14.7\%) & 1,794 & (12.0\%) & 360 & (11.4\%) \\
\quad HIV coinfection & 32,077 & (26.8\%) & 3,981 & (26.9\%) & 717 & (23.2\%) \\
\quad Retreatment & 8,351 & (7.0\%) & 1,262 & (8.7\%) & 290 & (9.4\%) \\
\hline\hline
\end{tabular}}
\end{table}

\renewcommand{\thetable}{1b}
\begin{table}[H]
    \centering
    \caption{Treatment outcomes across TIBU and paper registry cohorts.}
    \label{tbl:outcomes}
    \scriptsize{
\begin{tabular}{lrr|rr|rr}
\hline\hline \\[-8pt]
Characteristic & \multicolumn{4}{c}{Digital} & \multicolumn{2}{c}{Paper} \\
 & \multicolumn{2}{c}{\shortstack{All \\ (N=119,653)}} & \multicolumn{2}{c}{\shortstack{Study clinics \\ (N=14,962)}} & \multicolumn{2}{c}{\shortstack{Study clinics \\ (N=3,263)}} \\
 & $N$ & (\%) & $N$ & (\%) & $N$ & (\%) \\ \hline \\[-8pt]
\textit{Outcomes} & & & & & & \\
\quad Cured & 25,399 & (21.2\%) & 6,757 & (45.2\%) & 1,114 & (34.1\%) \\
\quad Treatment completed & 31,719 & (26.5\%) & 6,637 & (44.4\%) & 1,468 & (45.0\%) \\
\quad Died & 4,911 & (4.1\%) & 765 & (5.1\%) & 181 & (5.5\%) \\
\quad Treatment failed & 382 & (0.3\%) & 123 & (0.8\%) & 34 & (1.0\%) \\
\quad Lost to follow-up & 3,594 & (3.0\%) & 680 & (4.5\%) & 125 & (3.8\%) \\
\quad Not completed & 47,184 & (39.4\%) & 0 & (0.0\%) & 87 & (2.7\%) \\
\quad Transfer out & 6,382 & (5.3\%) & 0 & (0.0\%) & 138 & (4.2\%) \\
\quad Not TB & 0 & (0.0\%) & 0 & (0.0\%) & 0 & (0.0\%) \\
\quad Missing & 82 & (0.1\%) & 0 & (0.0\%) & 116 & (3.6\%) \\
\hline\hline
\end{tabular}}
\end{table}
\setcounter{table}{1}
\renewcommand{\thetable}{\arabic{table}}

% =============================================================================
\section{Data Quality Assurance: Outcome Discrepancies}
% =============================================================================

\subsection{Crosstab of TIBU vs.\ Paper Registry Outcomes}

Table~\ref{tbl:crosstab} shows the full cross-tabulation of treatment outcomes
recorded in TIBU against outcomes recorded in the paper registry for all
matched patients ($N = 3{,}370$).

\begin{table}[H]
    \centering
    \caption{Cross-tabulation of treatment outcomes in TIBU vs.\ paper registry.}
    \label{tbl:crosstab}
    \resizebox{\textwidth}{!}{\newsavebox{\DQACrosstabBox}
\savebox{\DQACrosstabBox}{\scriptsize{
\begin{tabular}{l|rr|rr|rr|rr|rr|rr|rr|rr|r}
\hline \hline \\[-8pt]
Outcome in TIBU & \multicolumn{16}{c}{Treatment Outcome in Paper Registry} & Total\\
& \multicolumn{2}{c|}{Blank} & \multicolumn{2}{c|}{C} & \multicolumn{2}{c|}{D} & \multicolumn{2}{c|}{F} & \multicolumn{2}{c|}{LTFU} & \multicolumn{2}{c|}{NC} & \multicolumn{2}{c|}{TC} & \multicolumn{2}{c|}{TO} & \\ 
& $N$ & (\%) & $N$ & (\%) & $N$ & (\%) & $N$ & (\%) & $N$ & (\%) & $N$ & (\%) & $N$ & (\%) & $N$ & (\%) & \\ \hline \\[-8pt]
Blank & 82 & (76.6\%) & \textcolor{blue}{14} & \textcolor{blue}{(13.1\%)} & \textcolor{blue}{1} & \textcolor{blue}{(0.9\%)} & \textcolor{blue}{0} & \textcolor{blue}{(0.0\%)} & \textcolor{blue}{2} & \textcolor{blue}{(1.9\%)} & 0 & (0.0\%) & \textcolor{blue}{6} & \textcolor{blue}{(5.6\%)} & \textcolor{blue}{2} & \textcolor{blue}{(1.9\%)} & 107 \\
C & 41 & (3.6\%) & 995 & (87.0\%) & \textcolor{red}{1} & \textcolor{red}{(0.1\%)} & \textcolor{red}{3} & \textcolor{red}{(0.3\%)} & \textcolor{red}{7} & \textcolor{red}{(0.6\%)} & 0 & (0.0\%) & 55 & (4.8\%) & \textcolor{red}{42} & \textcolor{red}{(3.7\%)} & 1,144 \\
D & 4 & (2.2\%) & \textcolor{green!60!black}{2} & \textcolor{green!60!black}{(1.1\%)} & 166 & (90.2\%) & 0 & (0.0\%) & 1 & (0.5\%) & 0 & (0.0\%) & \textcolor{green!60!black}{8} & \textcolor{green!60!black}{(4.3\%)} & \textcolor{green!60!black}{3} & \textcolor{green!60!black}{(1.6\%)} & 184 \\
F & 0 & (0.0\%) & \textcolor{green!60!black}{1} & \textcolor{green!60!black}{(2.9\%)} & 1 & (2.9\%) & 25 & (73.5\%) & 5 & (14.7\%) & 0 & (0.0\%) & \textcolor{green!60!black}{2} & \textcolor{green!60!black}{(5.9\%)} & \textcolor{green!60!black}{0} & \textcolor{green!60!black}{(0.0\%)} & 34 \\
LTFU & 3 & (2.3\%) & \textcolor{green!60!black}{2} & \textcolor{green!60!black}{(1.5\%)} & 1 & (0.8\%) & 0 & (0.0\%) & 113 & (86.3\%) & 0 & (0.0\%) & \textcolor{green!60!black}{6} & \textcolor{green!60!black}{(4.6\%)} & \textcolor{green!60!black}{6} & \textcolor{green!60!black}{(4.6\%)} & 131 \\
NC & 23 & (20.4\%) & \textcolor{blue}{34} & \textcolor{blue}{(30.1\%)} & \textcolor{blue}{6} & \textcolor{blue}{(5.3\%)} & \textcolor{blue}{0} & \textcolor{blue}{(0.0\%)} & \textcolor{blue}{4} & \textcolor{blue}{(3.5\%)} & 1 & (0.9\%) & \textcolor{blue}{32} & \textcolor{blue}{(28.3\%)} & \textcolor{blue}{13} & \textcolor{blue}{(11.5\%)} & 113 \\
TC & 30 & (2.0\%) & 76 & (5.0\%) & \textcolor{red}{7} & \textcolor{red}{(0.5\%)} & \textcolor{red}{0} & \textcolor{red}{(0.0\%)} & \textcolor{red}{12} & \textcolor{red}{(0.8\%)} & 0 & (0.0\%) & 1,333 & (87.7\%) & \textcolor{red}{62} & \textcolor{red}{(4.1\%)} & 1,520 \\
TO & 7 & (2.7\%) & \textcolor{orange!85!black}{45} & \textcolor{orange!85!black}{(17.0\%)} & \textcolor{orange!85!black}{5} & \textcolor{orange!85!black}{(1.9\%)} & \textcolor{orange!85!black}{5} & \textcolor{orange!85!black}{(1.9\%)} & \textcolor{orange!85!black}{14} & \textcolor{orange!85!black}{(5.3\%)} & 0 & (0.0\%) & \textcolor{orange!85!black}{63} & \textcolor{orange!85!black}{(23.9\%)} & 125 & (47.3\%) & 264 \\
\hline \\[-8pt]
Total & 190 &  & 1,169 &  & 188 &  & 33 &  & 158 &  & 1 &  & 1,505 &  & 253 &  & 3,497 \\
\hline \hline
\end{tabular}}}
\begin{minipage}{\wd\DQACrosstabBox}
\usebox{\DQACrosstabBox}
\par\vspace{4pt}\noindent
\textit{Note:} Percentages are row shares. Transfer-out (TO) is a documented outcome; a disagreeing TIBU value is a discrepancy named after what TIBU recorded. \\\textcolor{red}{Red} = False positive (TIBU: success; paper: failure or transfer). \\\textcolor{green!60!black}{Green} = False negative (TIBU: failure; paper: success or transfer). \\\textcolor{orange!85!black}{Orange} = False transfer (TIBU: transfer; paper: a real success/failure). \\\textcolor{blue}{Blue} = Missing in TIBU (TIBU: no outcome; paper: outcome present).
\end{minipage}}
\end{table}

\subsection{Error Rates by Clinic Characteristics}

Table~\ref{tbl:clinic_errors} reports false-negative, false-positive, and false-missing rates stratified by province, urbanicity, and clinic size.

\begin{table}[H]
    \centering
    \caption{DQA error rates by province, urbanicity, and clinic size.}
    \label{tbl:clinic_errors}
    \resizebox{\textwidth}{!}{\newsavebox{\ClinicErrBox}
\savebox{\ClinicErrBox}{\footnotesize{
\begin{tabular}{l|rr|rr|rr|rr|rr}
\hline\hline \\[-8pt]
& \multicolumn{2}{c|}{} & \multicolumn{2}{c|}{False positive} & \multicolumn{2}{c|}{False negative} & \multicolumn{2}{c|}{Missing in TIBU} & \multicolumn{2}{c}{False transfer} \\
& $N$ & (\%) & $N$ & (\%) & $N$ & (\%) & $N$ & (\%) & $N$ & (\%) \\
\hline \\[-8pt]
\textit{By province} &  &  &  &  &  &  &  &  &  & \\
\quad Central & 27 & (0.8\%) & 0 & (0.0\%) & 0 & (0.0\%) & 13 & (48.1\%) & 0 & (0.0\%) \\
\quad Coast & 708 & (21.4\%) & 18 & (2.5\%) & 4 & (0.6\%) & 13 & (1.8\%) & 6 & (0.8\%) \\
\quad Eastern North & 2 & (0.1\%) & 0 & (0.0\%) & 0 & (0.0\%) & 0 & (0.0\%) & 0 & (0.0\%) \\
\quad Eastern South & 264 & (8.0\%) & 8 & (3.0\%) & 2 & (0.8\%) & 24 & (9.1\%) & 12 & (4.5\%) \\
\quad Nairobi & 535 & (16.2\%) & 18 & (3.4\%) & 7 & (1.3\%) & 4 & (0.7\%) & 89 & (16.6\%) \\
\quad North Eastern & 821 & (24.8\%) & 66 & (8.0\%) & 6 & (0.7\%) & 2 & (0.2\%) & 19 & (2.3\%) \\
\quad Nyanza North & 392 & (11.9\%) & 6 & (1.5\%) & 7 & (1.8\%) & 3 & (0.8\%) & 2 & (0.5\%) \\
\quad Nyanza South & 40 & (1.2\%) & 1 & (2.5\%) & 0 & (0.0\%) & 0 & (0.0\%) & 0 & (0.0\%) \\
\quad Rift Valley North & 44 & (1.3\%) & 0 & (0.0\%) & 0 & (0.0\%) & 0 & (0.0\%) & 0 & (0.0\%) \\
\quad Rift Valley South & 425 & (12.9\%) & 15 & (3.5\%) & 3 & (0.7\%) & 37 & (8.7\%) & 4 & (0.9\%) \\
\quad Western & 31 & (0.9\%) & 2 & (6.5\%) & 1 & (3.2\%) & 1 & (3.2\%) & 0 & (0.0\%) \\
\rule{0pt}{14pt}\textit{By urban/rural} &  &  &  &  &  &  &  &  &  & \\
\quad Urban & 1,118 & (33.8\%) & 30 & (2.7\%) & 10 & (0.9\%) & 21 & (1.9\%) & 95 & (8.5\%) \\
\quad Rural & 2,167 & (65.5\%) & 104 & (4.8\%) & 20 & (0.9\%) & 76 & (3.5\%) & 37 & (1.7\%) \\
\rule{0pt}{14pt}\textit{By clinic size} &  &  &  &  &  &  &  &  &  & \\
\quad Large clinic & 1,273 & (38.5\%) & 39 & (3.1\%) & 10 & (0.8\%) & 24 & (1.9\%) & 90 & (7.1\%) \\
\quad Small clinic & 2,033 & (61.5\%) & 95 & (4.7\%) & 20 & (1.0\%) & 90 & (4.4\%) & 42 & (2.1\%) \\
\rule{0pt}{14pt}All & 3,306 & (100.0\%) & 134 & (4.1\%) & 30 & (0.9\%) & 114 & (3.4\%) & 132 & (4.0\%) \\
\hline\hline
\end{tabular}}}
\begin{minipage}{\wd\ClinicErrBox}
\usebox{\ClinicErrBox}
\par\vspace{4pt}\noindent
\footnotesize{\textit{Note:} Urban/rural classification is based on clinic records. Large clinics are those with at least 7 patients registered in TIBU (median); small clinics are those with fewer.}
\end{minipage}}
\end{table}

\subsection{Error Rates by Patient Characteristics}

Table~\ref{tbl:patient_errors} reports false-negative, false-positive, and
false-missing rates broken down by individual characteristics and disease
characteristics.

\begin{table}[H]
    \centering
    \caption{DQA error rates by individual and disease characteristics.}
    \label{tbl:patient_errors}
    \resizebox{\textwidth}{!}{\newsavebox{\PatientErrBox}
\savebox{\PatientErrBox}{\footnotesize{
\begin{tabular}{l|rr|rr|rr|rr||rr|rr|rr|rr}
\hline\hline \\[-8pt]
& \multicolumn{8}{c||}{\textbf{All clinics}} & \multicolumn{8}{c}{\textbf{Excluding clinics 0 and 115}} \\
\cline{2-9}\cline{10-17} \\[-8pt]
& \multicolumn{2}{c|}{} & \multicolumn{2}{c|}{False positive} & \multicolumn{2}{c|}{False negative} & \multicolumn{2}{c||}{Missing in TIBU} & \multicolumn{2}{c|}{} & \multicolumn{2}{c|}{False positive} & \multicolumn{2}{c|}{False negative} & \multicolumn{2}{c}{Missing in TIBU} \\
& $N$ & (\%) & $N$ & (\%) & $N$ & (\%) & $N$ & (\%) & $N$ & (\%) & $N$ & (\%) & $N$ & (\%) & $N$ & (\%) \\
\hline \\[-8pt]
\textit{Individual characteristics} &  &  &  &  &  &  &  &  &  &  &  &  &  &  &  & \\
\quad Male & 1,872 & (63.4\%) & 19 & (1.0\%) & 14 & (0.7\%) & 68 & (3.6\%) & 1,637 & (63.0\%) & 18 & (1.1\%) & 14 & (0.9\%) & 20 & (1.2\%) \\
\rule{0pt}{14pt}\textit{Age group} &  &  &  &  &  &  &  &  &  &  &  &  &  &  &  & \\
\quad $<$15 & 420 & (14.2\%) & 1 & (0.2\%) & 3 & (0.7\%) & 5 & (1.2\%) & 366 & (14.1\%) & 1 & (0.3\%) & 3 & (0.8\%) & 3 & (0.8\%) \\
\quad 15--34 & 1,334 & (45.2\%) & 18 & (1.3\%) & 11 & (0.8\%) & 35 & (2.6\%) & 1,170 & (45.0\%) & 17 & (1.5\%) & 11 & (0.9\%) & 11 & (0.9\%) \\
\quad 35--54 & 866 & (29.3\%) & 6 & (0.7\%) & 5 & (0.6\%) & 29 & (3.3\%) & 755 & (29.0\%) & 6 & (0.8\%) & 5 & (0.7\%) & 11 & (1.5\%) \\
\quad 55+ & 297 & (10.1\%) & 6 & (2.0\%) & 2 & (0.7\%) & 8 & (2.7\%) & 273 & (10.5\%) & 6 & (2.2\%) & 2 & (0.7\%) & 3 & (1.1\%) \\
\rule{0pt}{14pt}\textit{Disease characteristics} &  &  &  &  &  &  &  &  &  &  &  &  &  &  &  & \\
\quad Bacteriologically confirmed & 2,041 & (69.1\%) & 25 & (1.2\%) & 14 & (0.7\%) & 61 & (3.0\%) & 1,788 & (68.8\%) & 24 & (1.3\%) & 14 & (0.8\%) & 24 & (1.3\%) \\
\quad Extrapulmonary & 331 & (11.2\%) & 2 & (0.6\%) & 0 & (0.0\%) & 8 & (2.4\%) & 302 & (11.6\%) & 2 & (0.7\%) & 0 & (0.0\%) & 5 & (1.7\%) \\
\quad Retreatment & 275 & (9.3\%) & 5 & (1.8\%) & 1 & (0.4\%) & 5 & (1.8\%) & 250 & (9.6\%) & 5 & (2.0\%) & 1 & (0.4\%) & 2 & (0.8\%) \\
\quad HIV positive & 667 & (22.6\%) & 6 & (0.9\%) & 8 & (1.2\%) & 20 & (3.0\%) & 577 & (22.2\%) & 5 & (0.9\%) & 8 & (1.4\%) & 9 & (1.6\%) \\
\rule{0pt}{14pt}All & 2,953 & (100.0\%) & 31 & (1.0\%) & 21 & (0.7\%) & 112 & (3.8\%) & 2,599 & (100.0\%) & 30 & (1.2\%) & 21 & (0.8\%) & 62 & (2.4\%) \\
\hline\hline
\end{tabular}}}
\begin{minipage}{\wd\PatientErrBox}
\usebox{\PatientErrBox}
\par\vspace{4pt}\noindent
\footnotesize{\textit{Note:} Each row shows the count and rate of records matching the row label, and the count and rate of each error type within those records. Rows are not mutually exclusive (e.g.\ a male HIV-positive patient appears in both rows). The right panel excludes clinic IDs 0 and 115 (the two clinics responsible for the COVID-era missing-in-TIBU spike).}
\end{minipage}}
\end{table}

\subsection{Age of Record by Error Type}

Table~\ref{tbl:lag} reports the age of record (days from date of registration
to TIBU outcome date) stratified by error type.

\begin{table}[H]
    \centering
    \caption{Age of record by error category (in days from registration to TIBU outcome date).}
    \label{tbl:lag}
    \newsavebox{\CorrLagBox}
\savebox{\CorrLagBox}{\scriptsize{
\begin{tabular}{l|rr|rrr|rrr}
\hline\hline \\[-8pt]
 & & & \multicolumn{3}{c|}{Record age (days)}   & \multicolumn{3}{c}{Reporting lag (days)} \\
Error type & $N$ & (\%) & $n$ & Mean & Median & $n$ & Mean & Median \\
\hline \\[-8pt]
\textcolor{red}{False positive (Type I)} & 31 & (1.0\%) & 31 & 1020 & 982 & 31 & 873 & 807 \\
\textcolor{green!60!black}{False negative (Type II)} & 21 & (0.7\%) & 21 & 1033 & 1019 & 21 & 964 & 1005 \\
\textcolor{blue}{Missing in TIBU} & 112 & (3.8\%) & 112 & 920 & 844 & 110 & 749 & 690 \\
No error & 2,789 & (94.4\%) & 2,789 & 1029 & 1002 & 2,789 & 847 & 848 \\
\hline\hline
\end{tabular}}}
\begin{minipage}{\wd\CorrLagBox}
\usebox{\CorrLagBox}
\par\vspace{4pt}\noindent
\scriptsize{\textit{Note:} The TIBU snapshot is taken to be 1 January 2022, based on the latest treatment-outcome dates in the extract clustering in late December 2021. Record age is measured from registration to snapshot; reporting lag is measured from TIBU outcome date to snapshot. Column $n$ gives the number of records with a valid date; $N$ counts all records in the error category.}
\end{minipage}
\end{table}

\subsection{Error Rates by Age of Record}

\setcounter{figure}{1}
Figure~\ref{fig:age} plots DQA error rates by age of record (days from registration to TIBU outcome date), binned in 90-day (quarterly) intervals with a 3-quarter centered rolling mean ($N \geq 10$ per bin). Figure~\ref{fig:age_n100} restricts to bins with $N \geq 100$ to reduce noise from small samples.

\begin{figure}[H]
    \centering
    \includegraphics[width=0.85\textwidth]{fig_error_by_age.pdf}
    \caption{DQA error rates by age of record (bins with $N \geq 10$).}
    \label{fig:age}
\end{figure}

\begin{figure}[H]
    \centering
    \includegraphics[width=0.85\textwidth]{fig_error_by_age_n100.pdf}
    \caption{DQA error rates by age of record, restricted to bins with $N \geq 100$.}
    \label{fig:age_n100}
\end{figure}

\subsection{Error Rates Over Time}

Figure~\ref{fig:over_time} plots DQA error rates across the study period, binned by quarter with a 3-quarter centered rolling mean ($N \geq 10$ per quarter). Figure~\ref{fig:over_time_n100} restricts to quarters with $N \geq 100$ to reduce noise from small samples.

\begin{figure}[H]
    \centering
    \includegraphics[width=0.85\textwidth]{fig_error_over_time.pdf}
    \caption{DQA error rates over time, binned by quarter with a 3-quarter centered rolling mean ($N \geq 10$ per quarter).}
    \label{fig:over_time}
\end{figure}

\begin{figure}[H]
    \centering
    \includegraphics[width=0.85\textwidth]{fig_error_over_time_n100.pdf}
    \caption{DQA error rates over time, binned by quarter with a 3-quarter centered rolling mean ($N \geq 100$ per quarter).}
    \label{fig:over_time_n100}
\end{figure}

\end{document}
